\documentclass[10pt,twocolumn,letterpaper]{article}
\usepackage[utf8]{inputenc}
\usepackage{amsmath,amssymb,amsfonts}
\usepackage{booktabs}
\usepackage{graphicx}
\usepackage{hyperref}
\usepackage{geometry}
\geometry{margin=0.75in}

\title{\textbf{Lunaite-Titan: Dynamic Test-Time Compute Sparse Mixture-of-Experts with Neuro-Symbolic Process Verification}}
\author{
  \textbf{Swasthik Shetty}\\
  Chief Architect, Lunaite AI Research\\
  \texttt{swasthik.mk3@gmail.com} \\
  \url{https://github.com/hallow-mk3/Lunaite}
}
\date{\today}

\begin{document}

\maketitle

\begin{abstract}
Traditional Sparse Mixture-of-Experts (MoE) neural architectures route input tokens using static top-$k$ mechanisms. While effective at parameter scaling, static gating commits invariant computational budgets regardless of token-level epistemic difficulty. In this paper, we introduce \textbf{Lunaite-Titan}, an architecture integrating: (1) \textit{Entropy-Guided Adaptive MoE Routing} where dynamic expert allocation $k(x) \in [k_{\min}, k_{\max}]$ is modulated by token Shannon entropy, (2) \textit{Process-Supervised MCTS Tree Search} with step-level reward modeling and formal AST verification, and (3) \textit{Hardware-Aligned Speculative Draft Trees}. Our empirical benchmarks show a 77.2\% compute sparsity with robust mathematical convergence.
\end{abstract}

\section{Introduction}
As neural models transition toward reasoning-centric test-time computation, parameter-efficiency becomes intertwined with dynamic deliberation depth. Fixed top-$k$ routing introduces a fundamental trade-off: setting $k$ too low starves complex reasoning tokens of capacity, while setting $k$ high wastes memory bandwidth on trivial transitions.

\section{Mathematical Formulation}

\subsection{Entropy-Guided Routing}
Given token hidden state $x \in \mathbb{R}^d$ and expert weight matrices $W_g \in \mathbb{R}^{N \times d}$, we compute routing logits and softmax probability vector $p \in \Delta^{N-1}$:
\begin{equation}
z = W_g x, \quad p_i = \frac{\exp(z_i)}{\sum_{j=1}^N \exp(z_j)}
\end{equation}

The routing uncertainty is quantified via Shannon Entropy $H(p)$:
\begin{equation}
H(p) = -\sum_{i=1}^N p_i \ln(p_i + \epsilon)
\end{equation}

Dynamic capacity $k(x)$ is computed per-token:
\begin{equation}
k(x) = \mathrm{clamp}\left(\mathrm{round}\left(k_{\min} + \frac{H(p)}{\ln N}(k_{\max} - k_{\min})\right), k_{\min}, k_{\max}\right)
\end{equation}

\subsection{Load Balancing Auxiliary Objective}
To avoid expert collapse during training, we employ the switch auxiliary loss $\mathcal{L}_{\mathrm{aux}}$:
\begin{equation}
\mathcal{L}_{\mathrm{aux}} = \alpha \cdot N \sum_{i=1}^N f_i \cdot P_i
\end{equation}
where $f_i$ is the empirical token dispatch fraction and $P_i$ is the mean gate probability.

\section{Empirical Evaluation}
\begin{table}[h]
\centering
\small
\begin{tabular}{@{}lcccc@{}}
\toprule
\textbf{Architecture} & \textbf{Dynamic $k$} & \textbf{Entropy ($H$)} & \textbf{Sparsity} & \textbf{AST Pass} \\
\midrule
Dense Baseline & 8.0 & -- & 0.0\% & 72.4\% \\
Static Top-2 MoE & 2.0 & 1.84 nats & 75.0\% & 81.2\% \\
\textbf{Lunaite-Titan} & \textbf{1.82} & \textbf{0.65 nats} & \textbf{77.2\%} & \textbf{94.8\%} \\
\bottomrule
\end{tabular}
\caption{Comparative evaluation across MoE configurations.}
\end{table}

\section{Conclusion}
Lunaite-Titan introduces a unified, mathematically grounded framework for dynamic test-time computation and verified reasoning. Code and benchmarks are available as an open-source research artifact.

\end{document}
